\section{Specialization}
\label{section-specialization}

\begin{definition}
\label{definition-specialization}
Let $X$ be a topological space.
\begin{enumerate}
\item If $x, x' \in X$ then we say $x$ is a {\it specialization} of $x'$,
or $x'$ is a {\it generalization} of $x$ if $x \in \overline{\{x'\}}$.
Notation: $x' \leadsto x$.
\item A subset $T \subset X$ is {\it stable under specialization}
if for all $x' \in T$ and every specialization $x' \leadsto x$ we have
$x \in T$.
\item A subset $T \subset X$ is {\it stable under generalization}
if for all $x \in T$ and every generalization $x' \leadsto x$ we have
$x' \in T$.
\end{enumerate}
\end{definition}

\begin{lemma}
\label{lemma-open-closed-specialization}
Let $X$ be a topological space.
\begin{enumerate}
\item Any closed subset of $X$ is stable under specialization.
\item Any open subset of $X$ is stable under generalization.
\item A subset $T \subset X$ is stable under specialization
if and only if
the complement $T^c$ is stable under generalization.
\end{enumerate}
\end{lemma}

\begin{proof}
Let $F$ be a closed subset of $X$, if $y\in F$ then $\{y\} \subset F$, so
$\overline{\{y\}} \subset \overline{F} = F$ as $F$ is closed. Thus for all
specialization $x$ of $y$, we have $x\in F$.

\medskip\noindent
Let $x, y\in X$ such that $x\in \overline{\{y\}}$ and let $T$ be a subset of
$X$. Saying that $T$ is stable under specialization means that $y\in T$ implies
$x\in T$ and reciprocally saying that $T$ is stable under generalization means
that $x\in T$ implies $y\in T$. Therefore (3) is proven using contraposition.

\medskip\noindent
The second property follows from (1) and (3) by considering the complement.
\end{proof}

\begin{lemma}
\label{lemma-stable-specialization}
Let $T \subset X$ be a subset of a topological space $X$.
The following are equivalent
\begin{enumerate}
\item $T$ is stable under specialization, and
\item $T$ is a (directed) union of closed subsets of $X$.
\end{enumerate}
\end{lemma}

\begin{proof}
Suppose that $T$ is stable under specialization, then for all $y\in T$ we have
$\overline{\{y\}} \subset T$. Thus $T = \bigcup_{y\in T} \overline{\{y\}}$
which is an union of closed subsets of $X$. Reciprocally, suppose that $T =
\bigcup_{i\in I}F_i$ where $F_i$ are closed subsets of $X$. If $y\in T$ then
there exists $i\in I$ such that $y\in F_i$. As $F_i$ is closed, we have
$\overline{\{y\}} \subset F_i \subset T$, which proves that $T$ is stable under
specialization.
\end{proof}

\begin{definition}
\label{definition-lift-specializations}
Let $f : X \to Y$ be a continuous map of topological spaces.
\begin{enumerate}
\item We say that {\it specializations lift along $f$} or that $f$ is
{\it specializing} if given $y' \leadsto y$ in $Y$ and any $x'\in X$ with
$f(x') = y'$ there exists a specialization $x' \leadsto x$ of $x'$ in $X$ such
that $f(x) = y$.
\item We say that {\it generalizations lift along $f$} or that $f$ is
{\it generalizing} if given $y' \leadsto y$ in $Y$ and any $x\in X$ with
$f(x) = y$ there exists a generalization $x' \leadsto x$ of $x$ in $X$ such
that $f(x') = y'$.
\end{enumerate}
\end{definition}

\begin{lemma}
\label{lemma-lift-specialization-composition}
Suppose $f : X \to Y$ and $g : Y \to Z$ are continuous maps
of topological spaces. If specializations lift along both $f$ and $g$
then specializations lift along $g \circ f$. Similarly for
``generalizations lift along''.
\end{lemma}

\begin{proof}
Let $z'\leadsto z$ be a specialization in $Z$ and let $x' \in X$ such as
$g\circ f (x') = z'$. Then because specializations lift along $g$, there exists
a specialization $f(x') \leadsto y$ of $f(x')$ in $Y$ such that $g(y) = z$.
Likewise, because specializations lift along $f$, there exists a specialization
$x' \leadsto x$ of $x'$ in $X$ such that $f(x) = y$. It provides a
specialization $x' \leadsto x$ of $x'$ in $X$ such that $g\circ f(x) = z$. In
other words, specialization lift along $g\circ f$.
\end{proof}

\begin{lemma}
\label{lemma-lift-specializations-images}
Let $f : X \to Y$ be a continuous map of topological spaces.
\begin{enumerate}
\item If specializations lift along $f$, and if $T \subset X$
is stable under specialization, then $f(T) \subset Y$ is
stable under specialization.
\item If generalizations lift along $f$, and if $T \subset X$
is stable under generalization, then $f(T) \subset Y$ is
stable under generalization.
\end{enumerate}
\end{lemma}

\begin{proof}
Let $y' \leadsto y$ be a specialization in $Y$ where $y'\in f(T)$ and let
$x'\in T$ such that $f(x') = y'$. Because specialization lift along $f$, there
exists a specialization $x'\leadsto x$ of $x'$ in $X$ such that $f(x) = y$. But
$T$ is stable under specialization so $x\in T$ and then $y \in f(T)$. Therefore
$f(T)$ is stable under specialization.

\medskip\noindent
The proof of (2) is identical, using that generalizations lift along $f$.
\end{proof}

\begin{lemma}
\label{lemma-closed-open-map-specialization}
Let $f : X \to Y$ be a continuous map of topological spaces.
\begin{enumerate}
\item If $f$ is closed then specializations lift along $f$.
\item If $f$ is open, $X$ is a Noetherian topological space,
each irreducible closed subset of $X$ has a generic point,
and $Y$ is Kolmogorov then generalizations lift along $f$.
\end{enumerate}
\end{lemma}

\begin{proof}
Assume $f$ is closed. Let $y' \leadsto y$ in $Y$ and any $x'\in X$ with
$f(x') = y'$ be given. Consider the closed subset $T = \overline{\{x'\}}$
of $X$. Then $f(T) \subset Y$ is a closed subset, and $y' \in f(T)$.
Hence also $y \in f(T)$. Hence $y = f(x)$ with $x \in T$, i.e.,
$x' \leadsto x$.

\medskip\noindent
Assume $f$ is open, $X$ Noetherian, every irreducible closed subset of $X$
has a generic point, and $Y$ is Kolmogorov.
Let $y' \leadsto y$ in $Y$ and any $x \in X$ with
$f(x) = y$ be given. Consider $T = f^{-1}(\{y'\}) \subset X$.
Take an open neighbourhood $x \in U \subset X$ of $x$.
Then $f(U) \subset Y$ is open and $y \in f(U)$. Hence also $y' \in f(U)$.
In other words, $T \cap U \not = \emptyset$. This proves that
$x \in \overline{T}$. Since $X$ is Noetherian, $T$ is Noetherian
(Lemma \ref{lemma-Noetherian}).
Hence it has a decomposition $T = T_1 \cup \ldots \cup T_n$ into irreducible
components. Then correspondingly
$\overline{T} = \overline{T_1} \cup \ldots \cup \overline{T_n}$.
By the above $x \in \overline{T_i}$ for some $i$. By assumption
there exists a generic point $x' \in \overline{T_i}$, and
we see that $x' \leadsto x$. As $x' \in \overline{T}$ we see that
$f(x') \in \overline{\{y'\}}$. Note that
$f(\overline{T_i}) = f(\overline{\{x'\}}) \subset \overline{\{f(x')\}}$.
If $f(x') \not = y'$, then since $Y$ is Kolmogorov $f(x')$ is not a generic
point of the irreducible closed subset $\overline{\{y'\}}$ and the inclusion
$\overline{\{f(x')\}} \subset \overline{\{y'\}}$
is strict, i.e., $y' \not \in f(\overline{T_i})$.
This contradicts the fact that $f(T_i) = \{y'\}$.
Hence $f(x') = y'$ and we win.
\end{proof}

\begin{lemma}
\label{lemma-quotient-kolmogorov}
Suppose that $s, t : R \to U$ and $\pi : U \to X$ are continuous maps
of topological spaces such that
\begin{enumerate}
\item $\pi$ is open,
\item $U$ is sober,
\item $s, t$ have finite fibres,
\item generalizations lift along $s, t$,
\item $(t, s)(R) \subset U \times U$ is an equivalence relation on $U$ and
$X$ is the quotient of $U$ by this equivalence relation (as a set).
\end{enumerate}
Then $X$ is Kolmogorov.
\end{lemma}

\begin{proof}
Properties (3) and (5) imply that a point $x$ corresponds to an
finite equivalence class $\{u_1, \ldots, u_n\} \subset U$ of the equivalence
relation. Suppose that $x' \in X$ is a second point corresponding to
the equivalence class $\{u'_1, \ldots, u'_m\} \subset U$.
Suppose that $u_i \leadsto u'_j$ for some $i, j$. Then for any
$r' \in R$ with $s(r') = u'_j$ by (4) we can find $r \leadsto r'$
with $s(r) = u_i$. Hence $t(r) \leadsto t(r')$. Since
$\{u'_1, \ldots, u'_m\} = t(s^{-1}(\{u'_j\}))$ we conclude that
every element of $\{u'_1, \ldots, u'_m\}$ is the specialization of
an element of $\{u_1, \ldots, u_n\}$.
Thus $\overline{\{u_1\}} \cup \ldots \cup \overline{\{u_n\}}$ is
a union of equivalence classes, hence of the form $\pi^{-1}(Z)$
for some subset $Z \subset X$. By (1) we see that $Z$ is closed in $X$
and in fact $Z = \overline{\{x\}}$ because $\pi(\overline{\{u_i\}})
\subset \overline{\{x\}}$ for each $i$. In other words, $x \leadsto x'$
if and only if some lift of $x$ in $U$ specializes to some lift of
$x'$ in $U$, if and only if every lift of $x'$ in $U$ is a specialization
of some lift of $x$ in $U$.

\medskip\noindent
Suppose that both $x \leadsto x'$ and $x' \leadsto x$. Say $x$ corresponds
to $\{u_1, \ldots, u_n\}$ and $x'$ corresponds to $\{u'_1, \ldots, u'_m\}$
as above. Then, by the results of the preceding paragraph, we can find a
sequence
$$
\ldots \leadsto u'_{j_3} \leadsto u_{i_3} \leadsto u'_{j_2} \leadsto
u_{i_2} \leadsto u'_{j_1} \leadsto u_{i_1}
$$
which must repeat, hence by (2) we conclude that
$\{u_1, \ldots, u_n\} = \{u'_1, \ldots, u'_m\}$, i.e., $x = x'$.
Thus $X$ is Kolmogorov.
\end{proof}


\begin{lemma}
\label{lemma-dimension-specializations-lift}
Let $f : X \to Y$ be a morphism of topological spaces.
Suppose that $Y$ is a sober topological space, and $f$ is surjective.
If either specializations or generalizations lift along $f$, then
$\dim(X) \geq \dim(Y)$.
\end{lemma}

\begin{proof}
Assume specializations lift along $f$.
Let $Z_0 \subset Z_1 \subset \ldots Z_e \subset Y$ be a chain
of irreducible closed subsets of $X$. Let $\xi_e \in X$ be a point
mapping to the generic point of $Z_e$. By assumption there
exists a specialization
$\xi_e \leadsto \xi_{e - 1}$ in $X$ such that $\xi_{e - 1}$ maps to
the generic point of $Z_{e - 1}$. Continuing in this manner we find
a sequence of specializations
$$
\xi_e \leadsto \xi_{e - 1} \leadsto \ldots \leadsto \xi_0
$$
with $\xi_i$ mapping to the generic point of $Z_i$.
This clearly implies the sequence of irreducible closed
subsets
$$
\overline{\{\xi_0\}} \subset
\overline{\{\xi_1\}} \subset \ldots
\overline{\{\xi_e\}}
$$
is a chain of length $e$ in $X$.
The case when generalizations lift along $f$ is similar.
\end{proof}

\begin{lemma}
\label{lemma-characterize-closed-Noetherian}
Let $X$ be a Noetherian sober topological space.
Let $E \subset X$ be a subset of $X$.
\begin{enumerate}
\item If $E$ is constructible and stable under specialization, then
$E$ is closed.
\item If $E$ is constructible and stable under generalization, then
$E$ is open.
\end{enumerate}
\end{lemma}

\begin{proof}
Let $E$ be constructible and stable under generalization.
Let $Y \subset X$ be an irreducible closed subset with generic point
$\xi \in Y$. If $E \cap Y$ is nonempty, then it contains $\xi$ (by
stability under generalization) and hence is dense in $Y$, hence it
contains a nonempty open of $Y$, see
Lemma \ref{lemma-characterize-constructible-Noetherian}.
Thus $E$ is open by
Lemma \ref{lemma-characterize-open-Noetherian}.
This proves (2). To prove (1) apply (2) to the complement of $E$ in $X$.
\end{proof}
